\section*{Chapitre \ref{ch:dv:trees-and-finite-difference-pricers-in-practice} --- Arbres et pricers aux différences finies en pratique}

\begin{solution}{exo:dv:trees-and-finite-difference-pricers-in-practice:1}
Les nœuds terminaux de l'arbre se situent soit sur le strike, soit de part et d'autre de celui-ci selon la parité de $n$~: lorsque le strike coïncide avec un nœud, la rupture de pente
est représentée d'une façon, et lorsqu'elle tombe entre deux nœuds d'une autre, de sorte que l'erreur de discrétisation du payoff change de signe. L'erreur est d'ordre $1/n$ avec une
constante qui dépend de la parité.
\end{solution}

\begin{solution}{exo:dv:trees-and-finite-difference-pricers-in-practice:2}
$\log_2(0.30/0.091)=1.72$~; la division par deux suivante devrait donner environ $0.091/2^{1.72}=0.028$ cent (la grille donne 0.029).
\end{solution}

\begin{solution}{exo:dv:trees-and-finite-difference-pricers-in-practice:3}
Un dividende en numéraire décale chaque nœud du même montant, et non de la même proportion~: les chemins montée-puis-descente et descente-puis-montée ne se rejoignent donc plus, et
l'arbre cesse de se recombiner (sa taille croît exponentiellement). La grille est fixée en spot~: on applique $V(t_d^-,S)=V(t_d^+,S-D)$ par interpolation et
poursuit le calcul.
\end{solution}

\begin{solution}{exo:dv:trees-and-finite-difference-pricers-in-practice:4}
Avec $V_h=V+ch^2$~: $V\approx V_{1600}+(V_{1600}-V_{800})/3=6.090331+0.000068/3=6.0903535$.
\end{solution}

\begin{solution}{exo:dv:trees-and-finite-difference-pricers-in-practice:5}
La substitution arrière doit d'abord atteindre la région d'exercice, afin que chaque valeur y soit fixée par la contrainte avant que les valeurs de continuation voisines
ne soient calculées à partir d'elle~; pour un put, cette région se trouve à l'extrémité basse, de sorte que la substitution remonte et que l'élimination descend depuis le
sommet. Procéder dans l'autre sens reviendrait à calculer les valeurs de continuation à partir de voisins non projetés.
\end{solution}

\begin{solution}{exo:dv:trees-and-finite-difference-pricers-in-practice:6}
La sensibilité du prix au niveau de la barrière est d'ordre un (la dérivée de la valeur par rapport à $H$)~; une erreur d'une demi-cellule, $O(h)$,
sur le niveau se traduit donc par une erreur $O(h)$ sur le prix, qui domine l'erreur $O(h^2)$ du schéma.
\end{solution}

\begin{solution}{exo:dv:trees-and-finite-difference-pricers-in-practice:7}
Erreurs de 0.80, 0.35 et 0.16 cent pour 100, 200 et 400 nœuds~: ordre un (l'erreur est divisée par deux lorsque la grille double), contre 0.30, 0.091 et 0.029 avec
Brennan--Schwartz. Projeter après une résolution de Crank--Nicolson sans contrainte est incohérent avec le schéma.
\end{solution}

\begin{solution}{exo:dv:trees-and-finite-difference-pricers-in-practice:8}
La répétabilité n'est pas l'exactitude~: l'arbre est déterministe, il renvoie donc le même nombre à chaque exécution. À 2\,000 pas, son erreur est encore de l'ordre de 0.04 à
0.08 cent et alterne selon la parité du nombre de pas~; seule une étude de convergence contre une méthode plus précise permet de certifier les décimales.
\end{solution}

\begin{solution}{pb:dv:trees-and-finite-difference-pricers-in-practice:1}
\textbf{1.} 6.0864 et 6.0975~; la valeur est 6.0904.
\textbf{2.} $-0.40$ et $+0.72$ cent.
\textbf{3.} 1\,500 (et 1\,501).
\textbf{4.} $1500\times1501/2=1\,125\,750$.
\textbf{5.} La moyenne réduit le zigzag mais conserve l'erreur d'ordre un, dont la constante n'est pas connue à l'avance~; elle ne fournit aucune estimation d'erreur.
\textbf{6.} 1.04, 0.30, 0.091 et 0.029 cent.
\textbf{7.} 200 nœuds et pas~: 40\,600 mises à jour de nœuds.
\textbf{8.} Son erreur à 200 pas monte à 0.125 cent, au-delà de la tolérance~; et ses \omterm{def:dv:greeks-and-the-hedging-pnl:greeks}{grecques} au voisinage du strike oscillent.
\textbf{9.} 0.11 et 0.054 cent.
\textbf{10.} $V_{200}+(V_{200}-V_{100})/3$, à 0.021 cent près de la valeur.
\textbf{11.} 6.835 pour l'européenne, 7.407 pour l'américaine.
\textbf{12.} La barrière tombe entre deux nœuds~: erreurs de 0.16 à 200 nœuds contre 0.0002 lorsque la barrière est sur un nœud.
\textbf{13.} Gamma~: 0.217 sur le strike au lieu de 0.141 à une semaine de l'échéance, sans démarrage de Rannacher.
\textbf{14.} Une mise à jour d'un nœud de grille (une résolution tridiagonale) coûte quelques fois plus qu'une mise à jour d'un nœud d'arbre (une moyenne pondérée)~; le rapport dépend du langage et de
l'implémentation.
\textbf{15.} Elle réduit le temps de chaque méthode d'un facteur similaire, elle ne change donc pas la méthode gagnante~; elle déplace l'étude entière vers un domaine où un million de
mises à jour coûtent quelques millisecondes.
\textbf{16.} Aucune des deux~: elles s'écartent toutes deux d'environ un demi-cent, de part et d'autre~; il faut utiliser une méthode dont l'erreur est connue et faible, ou au moins procéder à une extrapolation.
\textbf{17.} Calculer le prix sur deux tailles de grille et prendre leur écart comme estimation de l'erreur~; comparer aux formes fermées lorsqu'elles existent~; placer le strike et
les barrières sur des nœuds~; consigner la grille utilisée.
\textbf{18.} Lorsque la simplicité prime sur le coût (enseignement, vérification rapide), et lorsque les règles d'exercice sont plus faciles à écrire nœud par nœud.
\textbf{19.} Arbre~: 1\,500 pas, 1\,125\,750 mises à jour de nœuds~; grille~: 200 nœuds et pas, 40\,600 mises à jour de nœuds~; la grille demande 28 fois moins de travail.
\textbf{20.} La méthode précise a une erreur connue et maîtrisée à la taille utilisée~; la méthode convergente promet seulement d'y parvenir à terme.
\end{solution}

\begin{solution}{iq:dv:trees-and-finite-difference-pricers-in-practice:1}
Crank--Nicolson amortit faiblement les composantes d'erreur haute fréquence~; un payoff à rupture de pente les injecte au strike, et près de l'échéance il y a trop peu de
pas pour qu'elles s'amortissent. Gamma, une dérivée seconde, les amplifie. Remède~: quelques demi-pas entièrement implicites (Rannacher) au début, et un
\omterm{def:dv:trees-and-finite-difference-pricers-in-practice:smoothing}{lissage du payoff}.
\iqlookfor{amortissement faible, la rupture de pente et les deux remèdes.}
\end{solution}

\begin{solution}{iq:dv:trees-and-finite-difference-pricers-in-practice:2}
Crank--Nicolson en $\ln S$ avec le strike sur un nœud, démarrage de Rannacher et \omterm{def:dv:trees-and-finite-difference-pricers-in-practice:bs}{algorithme de Brennan--Schwartz} pour le système tridiagonal sous contrainte (un
balayage par pas), vérifié sur deux tailles de grille~; pas une projection après une résolution sans contrainte, qui retombe à l'ordre un.
\iqlookfor{\omterm{def:dv:trees-and-finite-difference-pricers-in-practice:bs}{algorithme de Brennan--Schwartz} et détails de précision.}
\end{solution}

\begin{solution}{iq:dv:trees-and-finite-difference-pricers-in-practice:3}
Binomial~: ordre un, avec une oscillation entre nombres de pas pairs et impairs. Trinomial~: ordre un, plus lisse. Crank--Nicolson avec lissage, démarrage et
strike sur un nœud~: ordre deux. Pour un dixième de cent, la grille demande environ 30 fois moins de travail que l'arbre binomial dans le cas du chapitre.
\iqlookfor{les ordres et l'oscillation.}
\end{solution}

\begin{solution}{iq:dv:trees-and-finite-difference-pricers-in-practice:4}
Dividende~: une \omterm{def:dv:trees-and-finite-difference-pricers-in-practice:jump}{condition de saut} à la date ex, $V(t_d^-,S)=V(t_d^+,S-D)$, par interpolation, puis un redémarrage de Rannacher. Barrière~: une grille dont la frontière (ou un nœud) est
exactement sur la barrière, avec le rebate comme valeur au bord~; suivi discret de la barrière sous forme de projection aux dates.
\iqlookfor{\omterm{def:dv:trees-and-finite-difference-pricers-in-practice:jump}{condition de saut} et alignement.}
\end{solution}

\begin{solution}{iq:dv:trees-and-finite-difference-pricers-in-practice:5}
Études de convergence sur plusieurs tailles de grille, avec l'ordre observé comparé à la théorie~; accord avec les formes fermées et avec des méthodes indépendantes (un arbre, une
simulation de Monte Carlo)~; tests aux limites (volatilité nulle, strikes éloignés)~; régularité des \omterm{def:dv:greeks-and-the-hedging-pnl:greeks}{grecques}~; tests de non-régression calés sur des valeurs certifiées~; et les mêmes contrôles pour les
versions de production (C++) et de recherche (Python).
\iqlookfor{ordre de convergence, références indépendantes, non-régression.}
\end{solution}

\begin{solution}{iq:dv:trees-and-finite-difference-pricers-in-practice:6}
Pistes à examiner~: la \omterm{def:dv:trees-and-finite-difference-pricers-in-practice:jump}{condition de saut} du dividende appliquée au mauvais pas ou sans redémarrage de Rannacher (oscillations près de la frontière)~; interpolation à travers la
date ex, qui lisse la rupture de pente de la valeur d'exercice~; sens de Brennan--Schwartz non inversé pour un call (région d'exercice aux cours élevés)~; date ex absente de la
grille temporelle. Vérifier chacune sur une grille fine et sur un arbre avec le dividende.
\iqlookfor{traitement du dividende et sens d'élimination.}
\end{solution}
